S'ha preguntat a 200 persones si fan servir el transport públic habitualment ($T$) i si tenen cotxe ($C$).
En fan servir 120, i d'aquestes, 45 tenen cotxe. De les que no fan servir el transport públic, 60 tenen
cotxe. Es tria una persona a l'atzar.

\begin{apartats}

\apartat[1,25]{0,75}
Construeix la taula de contingència.

\begin{solucio}
\begin{center}
\begin{tabular}{l|cc|c}
 & Cotxe & Sense cotxe & Total\\ \hline
Transport públic & 45 & 75 & 120\\
No & 60 & 20 & 80\\ \hline
Total & 105 & 95 & 200
\end{tabular}
\end{center}
\end{solucio}

\begin{tria}{taula-tasca}
\itemtria{original}{1}{1,25}
Quina és la probabilitat que tingui cotxe? I si té cotxe, quina és la probabilitat que faci servir el
transport públic?

\begin{solucio}
$P(C)=\dfrac{105}{200}=0{,}525$.\\
$P(T\mid C)=\dfrac{45}{105}=\dfrac37\approx0{,}429$.
\end{solucio}

\itemtria{dos-titulars}{1}{1,25}
Un diari titula «El 43\,\% dels que tenen cotxe fan servir el transport públic», i un altre, «El 37,5\,\%
dels usuaris del transport públic tenen cotxe». Segons aquesta enquesta, té raó algun dels dos? O tots
dos?

\begin{solucio}
\textbf{Tots dos}: parlen de dues probabilitats condicionades diferents.\\
El primer és $P(T\mid C)=\dfrac{45}{105}\approx0{,}43$: el total de referència són els 105 que tenen
cotxe. El segon és $P(C\mid T)=\dfrac{45}{120}=0{,}375$: el total són els 120 usuaris del transport
públic. En general, $P(T\mid C)\neq P(C\mid T)$.
\end{solucio}
\end{tria}

\begin{nomesllarg}
\apartat{0,75}
Quina és la probabilitat que no faci servir el transport públic ni tingui cotxe? Són independents $T$ i
$C$?

\begin{solucio}
$P\left(\overline{T}\cap\overline{C}\right)=\dfrac{20}{200}=0{,}1$.\\
$P(T)\cdot P(C)=0{,}6\cdot0{,}525=0{,}315$, però $P(T\cap C)=\dfrac{45}{200}=0{,}225$: \textbf{no} són
independents.
\end{solucio}
\end{nomesllarg}

\end{apartats}
